\section{HW11: If-Based Flow Control (PL/NL/CEX)}

\subsection{Three styles for \cee{if-else} in ASM}
Given C \cee{if (cond) then; else else;}, three idiomatic ASM templates (Lctr 25):
\begin{tabular}{@{}l l p{0.40\columnwidth}@{}}
\toprule
Style & Branch on & Layout \\
\midrule
\textbf{PL} (pos. logic) & \textbf{same} CCS as C & else first, fall thru / B end / then \\
\textbf{NL} (neg. logic) & \textbf{inverse} CCS  & then first, fall thru / B end / else \\
\textbf{CEX} (cond exec) & no branch (IT block)  & all instrs in straight line \\
\bottomrule
\end{tabular}\\
NL reads closest to C (then-block falls through, matches \cee{if(c){...}}). PL is "branch when C-true". CEX is best when each arm is $\le$ 4 instructions total. CCS in \asm{IT..} = NL cond (the inverse of the C test).

\subsection{IT block syntax + hard rules}
\asm{IT\{x\{y\{z\}\}\} cond} -- up to \textbf{4} conditional instrs. First slot always \asm{T} (matches \asm{cond}); each later slot is \asm{T} (same cond) or \asm{E} (\textbf{inverse} cond).
\textbf{Common forms (T = same cond, E = inverse):} \asm{IT} (T), \asm{ITT} (T,T), \asm{ITE} (T,E), \asm{ITTT} (T,T,T), \asm{ITTE} (T,T,E), \asm{ITEE} (T,E,E), \asm{ITTEE} (T,T,E,E), \asm{ITTTT} (T,T,T,T).\\
\textbf{Hard rules (Lctr 24, p.\ 9):} max \textbf{4} cond instrs; \textbf{cond B} (\asm{Bcc}) only in \textbf{LAST} slot; \textbf{CBZ/CBNZ FORBIDDEN} inside IT (use \asm{cmp+Bcc}); each slot's CCS must match T/E role; cond LDR/STR (pre/post-idx) \textbf{legal} in IT.

\subsection{CBZ / CBNZ}
\asm{CBZ Rn,lbl} = "if Rn==0 goto lbl"; \asm{CBNZ Rn,lbl} = "if Rn!=0". \textbf{Forward only}, max \textbf{126 B} ahead. \asm{Rn} restricted to \textbf{R0--R7}. \textbf{Does NOT set flags}. Replaces \asm{cmp Rn,\#0 ; beq lbl} (1 instr vs 2). Forbidden inside IT. Use for null/loop guards.

\subsection{Sign-aware CCS (HW 11 always signed)}
Signed: \asm{GT/GE/LE/LT}; unsigned: \asm{HI/HS/LS/LO}. \textbf{Inverse pairs:} GT$\leftrightarrow$LE, GE$\leftrightarrow$LT, EQ$\leftrightarrow$NE. NL branches on the inverse. \textbf{Type drives 3 things:} CCS (GT/LE vs HI/LS), shift (ASR vs LSR), load (SB/SH vs B/H).

\subsection{Compound conditions}
\textbf{OR} (\cee{a||b}) -- short-circuit with TWO \asm{CMP}s, both branch to the SAME \asm{then} label; fall-through = else:
\begin{lstlisting}
    cmp r0, #20
    ble then         @ a true -> then
    cmp r0, #25
    bge then         @ b true -> then
    b   else_        @ neither -> else
\end{lstlisting}
\textbf{AND} (\cee{a\&\&b}) -- via De Morgan, branch to ELSE on inverse of either:
\begin{lstlisting}
    cmp r0, #LO
    blt else_        @ !a -> else (skip then)
    cmp r0, #HI
    bgt else_        @ !b -> else
    @ both true: fall through to then
\end{lstlisting}
For CEX of compound OR: 1st sub-test must be a plain \asm{Bcc} (cond B in middle of IT illegal); only the FINAL test can fan out into a full \asm{ITTEE}.

\subsection{Cond LDR/STR forms in IT}
All three addressing modes (\asm{[Rn,\#os]}, \asm{[Rn,\#os]!}, \asm{[Rn],\#os}) accept a CCS suffix (Lctr 24 Ex 3,4): e.g.\ \asm{ldrshgt r0,[r2,\#2]!}, \asm{ldrle r0,[r3],\#4}. Pre/post writeback honored only when cond true. Cond instr that fails: no write, no flag change, no exception.

\subsection{HW11 P1 -- \cee{x\_cmp\_y\_load}, three styles}
\begin{hwp}\textbf{P1.} \cee{int32\_t f(int32\_t x, int32\_t y, int16\_t *ptrA, int32\_t *ptrB)\{ if (x>y) return (int32\_t)*(++ptrA)*2; else return *(ptrB++)/2; \}}\\
AAPCS: R0=x, R1=y, R2=ptrA, R3=ptrB; ret in R0.\\
Then arm: \cee{*(++ptrA)} $\to$ \asm{ldrsh r0,[r2,\#2]!} (pre-idx, sign-ext halfword); \cee{*2} $\to$ \asm{lsl r0,r0,\#1}.\\
Else arm: \cee{*(ptrB++)} $\to$ \asm{ldr r0,[r3],\#4} (post-idx); signed \cee{/2} $\to$ \asm{asr r0,r0,\#1}.

\textbf{Part 1 -- PL (BGT, same CCS as C):}
\begin{lstlisting}
x_cmp_y_load_pl:
    cmp     r0, r1
    bgt     pl_then         @ x>y -> then
pl_else:                    @ else lays out FIRST
    ldr     r0, [r3], #4    @ *ptrB; ptrB++
    asr     r0, r0, #1      @ /2 signed
    b       pl_end
pl_then:
    ldrsh   r0, [r2, #2]!   @ ++ptrA; *ptrA (s-ext)
    lsl     r0, r0, #1      @ *2
pl_end:
    bx      lr
\end{lstlisting}

\textbf{Part 2 -- NL (BLE, inverse CCS, reads like C):}
\begin{lstlisting}
x_cmp_y_load_nl:
    cmp     r0, r1
    ble     nl_else         @ x<=y -> else (skip then)
nl_then:                    @ then falls through
    ldrsh   r0, [r2, #2]!
    lsl     r0, r0, #1
    b       nl_end
nl_else:
    ldr     r0, [r3], #4
    asr     r0, r0, #1
nl_end:
    bx      lr
\end{lstlisting}

\textbf{Part 3 -- CEX (\asm{ITTEE le}, no branches):}
\begin{lstlisting}
x_cmp_y_load_cex:
    cmp     r0, r1
    ittee   le              @ T=le,T=le,E=gt,E=gt
    ldrle   r0, [r3], #4    @ T1: else load (post-idx)
    asrle   r0, r0, #1      @ T2: else /2
    ldrshgt r0, [r2, #2]!   @ E1: then load (pre-idx)
    lslgt   r0, r0, #1      @ E2: then *2
    bx      lr
\end{lstlisting}
CCS in \asm{ITTEE} = \asm{le} (NL cond = inverse of C's \cee{>}). Each arm is 2 instrs $\Rightarrow$ \asm{ITTEE} (4 slots) is the natural fit. Cond pre-idx (\asm{ldrshgt!}) and post-idx (\asm{ldrle ,\#4}) both legal.
\end{hwp}

\subsection{HW11 P2 -- \cee{mathf} piecewise (3-way cascade)}
\begin{hwp}\textbf{P2.} \cee{f(x) = -10 if x<-5; +10 if x>=5; else 2x}.\\
\textbf{Part 1 (C):}
\begin{lstlisting}[language={}]
if (x < -5)      y = -10;
else if (x >= 5) y =  10;
else             y =  2*x;
\end{lstlisting}

\textbf{Part 2 -- CMP chain (NL at each arm):}
\begin{lstlisting}
mathf_cmp:
    cmp     r0, #-5         @ asm: cmn r0,#5
    bge     elseif          @ NL: x>=-5 skips 1st arm
    ldr     r0, =-10        @ x<-5 -> -10
    b       end
elseif:
    cmp     r0, #5
    blt     else_           @ NL: x<5 skips 2nd arm
    ldr     r0, =10         @ x>=5 -> +10
    b       end
else_:
    lsl     r0, r0, #1      @ 2*x
end:
    bx      lr
\end{lstlisting}
\textbf{Note:} \asm{cmp r0,\#-5} not directly encodable; assembler emits \asm{cmn r0,\#5} (flags = R0+5). Signed CCSs work unchanged.

\textbf{Part 3 -- CEX (chained \asm{ITT lt} + \asm{ITE ge}):}
\begin{lstlisting}
mathf_cex:
    cmp     r0, #-5
    itt     lt              @ T=lt, T=lt
    ldrlt   r1, =-10        @ T1: stash -10 in R1
    blt     end             @ T2: cond B *LAST slot only*
    @ fall-through: x >= -5
    cmp     r0, #5          @ fresh CMP, new flags
    ite     ge              @ T=ge, E=lt
    ldrge   r1, =10         @ T : x>=5, R1=10
    lsllt   r1, r0, #1      @ E : R1=2*x (uses ORIG x in R0)
end:
    mov     r0, r1          @ result -> ret reg
    bx      lr
\end{lstlisting}
\textbf{Key tricks:}
\begin{itemize}
\item \textbf{Stash result in R1} so R0 (=x) survives the 2nd CMP. \asm{lsllt r1,r0,\#1} reads original x.
\item \asm{blt end} is in \textbf{slot 2} of \asm{ITT} -- legal only because it's the \textbf{last} slot (Lctr 24 rule).
\item Without the cond B, fall-through after 1st IT would corrupt R1 in 2nd block.
\end{itemize}
\end{hwp}

\subsection{HW11 P3 -- compound OR \cee{x<=20 || x>=25}}
\begin{hwp}\textbf{P3.} \cee{int32\_t g(int x, uint16\_t *ptrA, uint32\_t *ptrB)\{ if (x<=20 || x>=25) return (int32\_t)*++ptrA*4; else return *++ptrB/4; \}}\\
AAPCS: R0=x, R1=ptrA, R2=ptrB; ret in R0.\\
Then arm: \cee{*++ptrA} $\to$ \asm{ldrh r0,[r1,\#2]!} (zero-ext); \cee{*4} $\to$ \asm{lsl r0,r0,\#2}.\\
Else arm: \cee{*++ptrB} $\to$ \asm{ldr r0,[r2,\#4]!}; UNSIGNED \cee{/4} $\to$ \asm{lsr r0,r0,\#2} (NOT \asm{ASR} -- ptrB is \cee{uint32\_t*}).

\textbf{Part 1 -- CMP short-circuit OR (both tests $\to$ same \cee{then}):}
\begin{lstlisting}
load_based_on_x_cmp:
    cmp     r0, #20
    ble     then            @ 1st OR true -> then
    cmp     r0, #25
    bge     then            @ 2nd OR true -> then
    b       else_           @ both false -> else
then:
    ldrh    r0, [r1, #2]!   @ ++ptrA; *ptrA (zero-ext)
    lsl     r0, r0, #2      @ *4
    b       end
else_:
    ldr     r0, [r2, #4]!   @ ++ptrB; *ptrB
    lsr     r0, r0, #2      @ /4 unsigned
end:
    bx      lr
\end{lstlisting}

\textbf{Part 2 -- CEX official (\asm{ITTT le} then \asm{ITTEE ge}, max-CEX):}
\begin{lstlisting}
load_based_on_x_cex:
    cmp     r0, #20
    ittt    le              @ T,T,T cond=le (x<=20 true)
    ldrhle  r0, [r1, #2]!   @ T1: then load (zero-ext)
    lslle   r0, r0, #2      @ T2: *4
    ble     end             @ T3: last-slot cond B (legal)
    @ fall-through: x>20 -> 2nd OR test
    cmp     r0, #25
    ittee   ge              @ T=ge,T=ge,E=lt,E=lt
    ldrhge  r0, [r1, #2]!   @ T: then load
    lslge   r0, r0, #2      @ T: *4
    ldrlt   r0, [r2, #4]!   @ E: else load
    lsrlt   r0, r0, #2      @ E: /4
end:
    bx      lr
\end{lstlisting}
\textbf{Why \asm{ITTT le}:} 1st OR test packs into IT, ending with cond \asm{B end} in slot 3 (legal only as LAST slot). 2nd test fans into \asm{ITTEE}. \emph{Repetition:} then-arm appears twice (once per IT block) -- official solution notes this is undesirable but inherent to packing both OR branches into CEX.

\textbf{Trace:}
\begin{tabular}{@{}c l l@{}}
\toprule
x & path & arm \\
\midrule
15 & \asm{bgt} N (15$\le$20) & inline then \\
20 & \asm{bgt} N            & inline then \\
21 & \asm{bgt} Y, ge N      & ITTEE E (else) \\
24 & \asm{bgt} Y, ge N      & ITTEE E (else) \\
25 & \asm{bgt} Y, ge Y      & ITTEE T (then) \\
30 & \asm{bgt} Y, ge Y      & ITTEE T (then) \\
\bottomrule
\end{tabular}
\end{hwp}

\subsection{Picker + pitfalls}
\textbf{Pick CEX} when both arms $\le$ 4 instrs total; \textbf{PL/NL} when arms larger or need general branch; \textbf{CMP short-circuit} for OR; \textbf{CMP+De Morgan} for AND (branch to else on $!a$/$!b$); \textbf{CBZ/CBNZ} for null/loop guards (forward $\le$126B, R0-R7, no flag effect).
\textbf{Gotchas:} (1) HI/LS on signed (or GT/LE on unsigned) miscompares near sign bit -- match CCS to type. (2) LSR on negative = huge positive; use ASR for signed div. (3) LDRH on \cee{int16\_t} loses sign -- use LDRSH. (4) Cond B in non-last IT slot = assembler err. (5) CBZ inside IT = forbidden. (6) In CEX cascade, stash result in R1 so R0 (=x) survives next CMP. (7) \asm{cmp Rn,\#-imm} rewritten to \asm{cmn Rn,\#imm} (signed CCS unchanged). (8) Pre/post-idx LDR both legal in IT, writeback honored only when cond true.

\subsection{One-line summary cards}
\begin{itemize}
\item \textbf{PL:} \asm{cmp; B<C-cond> then; <else>; b end; then: <then>; end:}
\item \textbf{NL:} \asm{cmp; B<inv-cond> else; <then>; b end; else: <else>; end:}
\item \textbf{CEX:} \asm{cmp; ITTEE <inv-cond>; <else x2 with cond>; <then x2 with inv-cond>}
\item \textbf{CEX cascade:} \asm{cmp; ITT lt; <stash-T>; blt end; cmp; ITE ge; <stash-T>; <stash-E>; end: mov r0,r1}
\item \textbf{OR (CMP):} \asm{cmp; ble then; cmp; bge then; b else}
\item \textbf{OR (CEX):} 1st = plain \asm{Bcc} to \asm{checkN}, 2nd packs into \asm{ITTEE}.
\end{itemize}
